\section*{Resumo}
\addcontentsline{toc}{section}{Resumo}
Este trabalho implementa e avalia um autopiloto didático em ESP32 com FreeRTOS, composto por fusão sensorial periódica, controle de atitude, navegação e segurança. Sensores touch geram comandos de rota, telemetria, estresse e emergência, enquanto sensores inerciais e motores são simulados. Foram analisados 18 ensaios em um núcleo, com duração de 30,000 a 35,015~s, abrangendo prioridades Deadline Monotonic (DM), customizadas e uma adaptação Rate Monotonic (RM), configurações preemptiva e cooperativa, frequências de 80, 160 e 240~MHz e dois ensaios adicionais sem time slicing. O processamento de 488 linhas de métricas permitiu comparar perdas de deadline, resposta, latência de ativação, jitter e ocupação aproximada das tarefas. CUSTOM e DM preemptivos a 160 e 240~MHz não registraram perdas nesta bateria. A 80~MHz, todas as condições avaliadas apresentaram perdas de navegação, e DM/RM preemptivos registraram atraso na conclusão do job de segurança. Os ensaios cooperativos registraram perdas da fusão inclusive a 240~MHz. A interpretação considera as limitações da referência temporal da fusão, da medição de duração com interferência e da quantidade de eventos de emergência. Os resultados evidenciam o compromisso entre prioridade da segurança, proteção do controle e capacidade de processamento, sem constituir garantia formal de pior caso.

\noindent\textbf{Palavras-chave:} sistemas em tempo real; ESP32; FreeRTOS; escalonamento; deadline; instrumentação temporal.
\clearpage
